\documentclass[12pt]{article}
\usepackage[margin=1in]{geometry}
\usepackage{amsmath}
\usepackage{fancyhdr}
\fancypagestyle{plain}{%
  \fancyhf{}
  \fancyhead[L]{POLISCI 450A}
  \fancyhead[R]{Fall 2026}
  %\fancyfoot[C]{\thepage}
}

\usepackage{titlesec}
\titleformat{\section}{\large\bfseries}{\thesection}{1em}{}
\usepackage{titling}
\setlength{\droptitle}{-3em}   % move title up


\title{\Large Lecture 2-1: Point Estimation}
\author{}
\date{}
\begin{document}
\maketitle
\vspace{-6em}

\noindent \textbf{Goal:} Introduce point estimation and key properties of estimators.\medskip

\begin{itemize}\setlength{\itemsep}{1em}
\item What is point estimation?
\begin{itemize}
\item Estimands, estimators, estimates
\item Statistical inference
\item The role of sampling
\end{itemize}
\item Properties of estimators
\begin{itemize}
\item Sampling distributions
\item Unbiasedness
\item Efficiency
\item Mean squared error
\end{itemize}
\item Evaluating estimators
\begin{itemize}
\item Analytic derivations
\item Simulation methods
\end{itemize}
\item Asymptotic properties
\begin{itemize}
\item Consistency
\item Convergence in probability
\item Law of large numbers
\end{itemize}
\item Asymptotic normality
\begin{itemize}
\item Central limit theorem
\item Approximate distributions
\item Simulation checks
\end{itemize}
\item Analogy principle
\begin{itemize}
\item Estimating population characteristics with sample counterparts
\end{itemize}
\end{itemize}

% End outline
\clearpage


\section*{Worksheet}


\begin{enumerate}\itemsep1em
\item Which of $\hat\mu_{1}, \hat\mu_{2}, \hat\mu_{3}, \hat\mu_{4}$ are unbiased? Show your work.
\vspace{20em}

\item Derive the variances of $\hat\mu_{1}, \hat\mu_{2}, \hat\mu_{3}, \hat\mu_{4}$. Show your work. What do you conclude?
\vspace{20em}
\clearpage 

\item Prove: $MSE(\hat{\theta}) = Var(\hat{\theta}) + \text{Bias}(\hat{\theta})^2.$
\vspace{25em}


\item Derive $MSE(\hat\mu_{1}), MSE(\hat\mu_{2}), MSE(\hat\mu_{3}), MSE(\hat\mu_{4})$
\vspace{20em}


\end{enumerate}

% \clearpage

% \section*{T/F}

% \begin{enumerate}\itemsep3em
% \item An estimand is a feature of the population distribution (e.g., $\mu$), whereas an estimator is a function of the data.
% \item If $Y_1,\dots,Y_n$ are i.i.d. with $E[Y_i]=\mu$, then $\overline{Y}_n$ is an unbiased estimator of $\mu$.
% \item The estimator that always returns the constant $4$ has zero variance and is therefore potentially a good estimator for the population mean $\mu$.
% \item Among $\widehat{\mu}_1=Y_1$ and $\widehat{\mu}_3=\overline{Y}_n$, both are unbiased, but $\overline{Y}_n$ is more efficient for $n>1$.
% \item $\widetilde{Y}_n=\frac{1}{n+5}\sum_{i=1}^n Y_i$ is biased for $\mu$ but can still be consistent for $\mu$.
% \item Unbiasedness implies consistency for any estimator.
% \item An estimator can be unbiased in every finite sample but still inconsistent in large samples.
% \item By the Weak Law of Large Numbers, $\overline{X}_n \xrightarrow{p} \mu$ for i.i.d. data with finite mean $\mu$, regardless of the population’s shape.
% \item The Central Limit Theorem states that $\sqrt{n}(\overline{X}_n-\mu)\xrightarrow{d}\mathcal{N}(0,\sigma^2)$ for i.i.d. data with finite variance.
% \item If an estimator is asymptotically normal, then for sufficiently large $n$ (e.g., $n = 1,000$), its sampling distribution is exactly normal.
% \end{enumerate}


\end{document}
